\documentclass{article}
\pdfpagewidth=8.5in
\pdfpageheight=11in

\usepackage{ijcai26}
\usepackage{times}
\usepackage{url}
\usepackage[hidelinks]{hyperref}
\usepackage[utf8]{inputenc}
\usepackage{graphicx}
\usepackage{amsmath}
\usepackage{booktabs}
\usepackage[brazilian]{babel}

\urlstyle{same}

% Cabeçalho do resumo em português
\renewenvironment{abstract}{\centerline{\Large\bf
Resumo}\vspace{0.5ex}\begin{quote}}{\par\end{quote}\vskip 0.5ex}

% Compatibilidade de metadados PDF
\ifdefined\pdfinfo
\pdfinfo{
/TemplateVersion (IJCAI.2026.0)
}
\fi

\title{Arquitetura e Avaliação de um Resolvedor para ARC-AGI-1 com LLMs: Fim a Fim e Geração de Código com Laço de Correção\\(Nível 2)}

\author{
    Mike Carvalho Borges
    \affiliations
    Especialização em Inteligência Artificial Avançada, UFRGS
    \emails
    mike.borges@ufrgs.br
}

\begin{document}

\maketitle

\begin{abstract}
Este trabalho apresenta o projeto, a implementação e a avaliação empírica de um resolvedor didático para o benchmark ARC-AGI-1, em Python e gerenciado via \texttt{uv}. O sistema emprega o Gemini 3.5 Flash Lite com raciocínio estendido (\textit{high thinking}) e implementa duas abordagens autônomas: inferência matricial direta (\textit{Fim a Fim}) e síntese de programas em Python (\textit{Geração de Código}) em \textit{sandbox} isolado com controle de tempo limite. O módulo de código integra um laço de correção reflexivo baseado em diagnósticos de erro nos pares de treino. Em 20 tarefas representativas, o resolvedor obteve 85\% de acurácia (17/20) em cada abordagem (90\% na união). O laço de correção recuperou 3 tarefas inicialmente incorretas (média de 1,40 tentativas/tarefa), e o ambiente identificou 2 casos de sobreajuste indutivo que acertaram 100\% do treino e falharam no teste.
\end{abstract}

\section{Introdução}

O ARC-AGI-1 (\textit{Abstraction and Reasoning Corpus}) foi proposto por François Chollet em 2019 como benchmark de inteligência fluida e eficiência amostral. Ao contrário de tarefas de classificação supervisionada, o ARC exige induzir regras de transformação espacial a partir de apenas 2 a 5 pares de treino de matrizes coloridas (valores 0 a 9) e aplicá-las a uma grade de teste inédita.

Para analisar a capacidade de modelos de fronteira com raciocínio estendido, construímos um resolvedor modular em torno do Gemini 3.5 Flash Lite (\texttt{gemini-3.5-flash-lite}, \textit{thinking high}), comparando duas estratégias:
\begin{enumerate}
    \item \textbf{Fim a Fim (\textit{End-to-End})}: Inferência direta em que o LLM raciocina e emite a matriz de teste serializada em JSON.
    \item \textbf{Geração de Código (\textit{Code Generation})}: Síntese de uma função pura \texttt{transform(grid)} executada em \textit{sandbox} e aprimorada via laço iterativo de correção reflexiva.
\end{enumerate}

\section{Arquitetura e Método do Resolvedor}

O resolvedor foi estruturado em um único ponto de acesso (\texttt{solver.py}), sob o gerenciador ultrarrápido \texttt{uv}.

\subsection{Fluxo de Dados e Representação Matricial}
O resolvedor lê o arquivo \texttt{tarefas.txt} e processa as tarefas em JSON. Cada grade é formatada com dimensões explícitas ($L \times C$), formatação visual tabular e legenda de cores (0: Preto a 9: Marrom), facilitando o raciocínio geométrico e topológico do modelo.

\subsection{Módulo Fim a Fim}
No modo Fim a Fim, o modelo deduz as invariâncias e retorna um objeto JSON com a chave \texttt{"output"}. O analisador sintático combina busca por blocos \texttt{```json}, inspeção de dicionários e expressões regulares de salvaguarda.

\subsection{Módulo de Código e Sandbox Seguro}
No modo de código, o LLM sintetiza a função pura \texttt{def transform(grid: list[list[int]]) -> list[list[int]]:}. A função é executada em processo isolado (\textit{subprocess}) com:
\begin{itemize}
    \item \textbf{Tempo Limite (\textit{Timeout})}: Interrupção forçada após 5,0s por grade para neutralizar laços infinitos.
    \item \textbf{Tratamento de Exceções}: Captura segura de \texttt{IndexError}, erros de recursão e tipos sem interromper o resolvedor.
\end{itemize}

\subsection{Laço de Correção Reflexivo}
O diferencial do método é o laço de correção (*reflection loop*). Quando o código sintetizado falha em exemplos de treino:
\begin{enumerate}
    \item O \textit{sandbox} intercepta saídas incorretas ou erros por par de treino.
    \item Constrói-se um prompt de diagnóstico com o código prévio, os contraexemplos (grade esperada vs. obtida) e a lista dos pares que passaram (preservando funcionalidades válidas).
    \item O modelo reflete sobre a falha e gera uma versão corrigida (limite configurável de até 3 tentativas).
    \item Ao atingir 100\% de acerto no treino, o laço é encerrado e a função é aplicada ao teste.
\end{enumerate}

\subsection{Protocolo Experimental}
O resolvedor respeita 1,0s de pausa entre requisições à API e audita tokens de entrada, raciocínio interno, saída, latência da API, tempo de \textit{sandbox}, número de tentativas e nota binária estrita (1,0 ou 0,0).

\begin{table}[t]
  \centering
  \scriptsize
  \setlength{\tabcolsep}{2.0pt}
  \begin{tabular}{lrrrrrr}
    \toprule
    Abordagem & Acurácia & Laço & Tent. & Tr/Te & Racioc. & Tempo (s) \\
    \midrule
    Fim a fim & 17/20 (85\%) & -- & 1,00 & 0 & 157.963 & 515,7 \\
    Código    & 17/20 (85\%) & 3  & 1,40 & 2 & 370.559 & 1.464,6 \\
    \bottomrule
  \end{tabular}
  \caption{Métricas consolidadas nas 20 tarefas (Gemini 3.5 Flash Lite, Raciocínio Alto). Laço: tarefas resolvidas c/ laço. Tent.: média de tentativas/tarefa. Tr/Te: casos com 100\% no treino e falha no teste.}
  \label{tab:resultados}
\end{table}

\section{Resultados e Discussão}

A Tabela~\ref{tab:resultados} consolida o desempenho obtido nas 20 tarefas avaliadas (10 de treino e 10 de avaliação). Ambas as abordagens atingiram \textbf{85,0\%} de acurácia individual (17/20), com 100\% de acerto (10/10) nas tarefas de treino (\texttt{training/}).

\subsection{Eficácia do Laço de Correção}
O laço de correção foi fundamental para o módulo de código. Em \textbf{3 tarefas}, o código inicial falhou no treino, mas o laço diagnosticou as inconsistências e encontrou a solução correta no teste:
\begin{itemize}
    \item \texttt{training/44d8ac46.json}: convergiu na 3ª tentativa após retificar a lógica de expansão geométrica.
    \item \texttt{training/0a938d79.json}: convergiu na 2ª tentativa, corrigindo o mapeamento de cores (4/4 no treino e nota 1,0 no teste).
    \item \texttt{evaluation/7c9b52a0.json}: convergiu na 3ª tentativa, ajustando o recorte dimensional e obtendo nota 1,0 no teste inédito.
\end{itemize}
A média foi de \textbf{1,40 tentativas por tarefa}. O laço elevou o consumo de tokens de raciocínio para 370.559 (aumento de 134,6\% frente ao modo fim a fim) e o tempo para 1.464,6s, refletindo o custo deliberativo adicional sobre contraexemplos.

\subsection{Sobreajuste ao Treino e Falha no Teste (Tr/Te)}
O resolvedor identificou \textbf{2 casos} em que o código alcançou 100\% de acerto no treino (3/3), mas errou o teste (\textit{Train Pass / Test Fail}):
\begin{itemize}
    \item \texttt{evaluation/d47aa2ff.json}: o código ajustou-se perfeitamente aos 3 pares de treino na 1ª tentativa, mas induziu critério espúrio que falhou no teste.
    \item \texttt{evaluation/f5aa3634.json}: a 1ª tentativa falhou (0/3); o laço corrigiu o código na 2ª tentativa (3/3 no treino), mas a regra ainda continha viés incompleto para o teste.
\end{itemize}
Esse resultado expõe a \textit{lacuna de generalização indutiva} em aprendizado com poucos exemplos (*few-shot*), onde múltiplas hipóteses satisfazem o treino mas apenas uma reflete o teste.

\subsection{Complementaridade e União dos Módulos}
Houve clara complementaridade entre as abordagens:
\begin{itemize}
    \item \texttt{evaluation/aab50785.json}: resolvida somente por Código (1,0). A tarefa requer extrair o interior de blocos $2\times2$; o fim a fim falhou por desalinhamento, enquanto o código calculou as coordenadas exatas.
    \item \texttt{evaluation/93b4f4b3.json}: resolvida somente por Fim a Fim (1,0). A tarefa envolve reflexão e simetria global, intuitiva visualmente para o LLM, mas difícil de parametrizar em código (esgotou as 3 tentativas).
\end{itemize}
A união das abordagens alcançou \textbf{18/20 tarefas resolvidas (90,0\%)}, evidenciando os ganhos de um sistema híbrido.

\section{Conclusão}

O resolvedor valida o potencial do Gemini 3.5 Flash Lite com raciocínio estendido no ARC-AGI-1. O laço de reflexão recuperou tarefas críticas ao inspecionar diagnósticos no \textit{sandbox}, ao mesmo tempo em que revelou casos de sobreajuste aos exemplos de demonstração. A combinação entre síntese determinística de código e intuição espacial direta constitui o caminho mais promissor para raciocínio abstrato robusto.

\paragraph{Uso de Inteligência Artificial:} O assistente de IA foi utilizado como par de programação para a arquitetura modular do resolvedor, laço de reflexão, \textit{sandbox} e automação experimental com \texttt{uv}.

\footnotesize
\begin{thebibliography}{99}
\setlength{\itemsep}{0pt}
\setlength{\parskip}{0pt}
\bibitem{chollet2019}
F.~Chollet.
\newblock On the measure of intelligence.
\newblock {\em arXiv:1911.01547}, 2019.

\bibitem{geminiteam2025}
{Gemini Team}.
\newblock Gemini: A family of highly capable multimodal models.
\newblock {\em arXiv:2312.11805}, 2025.
\end{thebibliography}

\end{document}
